\documentclass[11pt]{article}
\usepackage[margin=1in]{geometry}
\usepackage{amsmath}
\usepackage{booktabs}
\usepackage{hyperref}
\usepackage{parskip}
\usepackage{microtype}
\usepackage{xcolor}
\usepackage{listings}

\lstset{
  basicstyle=\ttfamily\small,
  backgroundcolor=\color{gray!10},
  frame=single,
  framerule=0.4pt,
  breaklines=true,
}

\title{Session Progress --- 2026-04-27}
\author{C.~elegans Locomotion Model}
\date{April 27, 2026}

\begin{document}
\maketitle
\tableofcontents
\vspace{1em}

\section*{Overview}
This session focused on building up oscillatory circuit complexity from a single neuron toward a single-segment network test, using a slow recovery variable $w$ as the primary mechanism for plateau termination.

\section{Single-Neuron Recovery Variable}

\textbf{Objective:} Confirm that the recovery variable $w$ reliably returns the neuron to rest after a plateau, even for weak (barely-threshold) stimuli.

\subsection*{Validated Parameters}
\begin{itemize}
  \item $\beta = 0.2\ \text{nS/mV}$ --- slope of $w_\infty(V) = \beta \cdot \max(V - T,\, 0)$ above threshold
  \item $\tau_w = 200\ \text{ms}$ --- recovery timescale
\end{itemize}

\subsection*{Key Finding}
Recovery is always mathematically guaranteed once the neuron reaches the plateau ($H$), because $f(H) = 0$ and the $-w/C$ term always pulls $V$ downward. The bottleneck is speed, not recoverability. There is an unstable fixed point in the $(V, w)$ phase plane at:
\begin{align}
  V^* &= H + \frac{\beta}{g \cdot m_3} = -41\ \text{mV} \\
  w^* &= \beta\,(V^* - T) = 0.8\ \text{pA}
\end{align}
The neuron returns to rest once $w$ accumulates past $w^* \approx 0.8\ \text{pA}$, which takes approximately 108 ms from $w = 0$ with these parameters.

\textbf{Test file:} \texttt{tests/test\_single\_neuron.py}

\section{Mutual Inhibition Circuit (II)}

\textbf{Objective:} Test whether $w$ alone can drive oscillations in a two-cell mutually inhibitory circuit.

\textbf{Result:} Failed. $w_\infty(H) = \beta\,(H - T) = 2.0\ \text{pA}$ saturates at or below the constant $I_\text{app}$ values tested ($2$--$4\ \text{pA}$). $w$ can never overcome a sustained constant drive, so neurons lock permanently to the plateau.

\textbf{Conclusion:} Synaptic fatigue (SF) remains necessary in the II circuit. SF was restored to \texttt{tests/test\_two\_II.py}. The failure is specific to constant $I_\text{app}$ --- in the full network, neurons receive transient synaptic input rather than a sustained DC drive, so the constraint is less severe.

\section{Excitatory-Inhibitory Circuit (EI)}

\textbf{Objective:} Test whether $w$ alone (no SF) can generate intrinsic oscillations in a two-cell EI circuit.

\textbf{Circuit:} $N_0$ (excitatory) $\to$ $N_1$ (inhibitory) $\to$ $N_0$. Constant $I_\text{app}$ on $N_0$ only; $N_1$ receives no direct input.

\textbf{Why the EI circuit works without SF:} $N_1$ has no constant drive, so $w$ does not need to overcome a sustained current --- it only needs to terminate $N_1$'s plateau, which it does.

\subsection*{Validated Parameters}
\begin{itemize}
  \item $G_\text{syne} = 0.063\ \text{nS}$, \quad $G_\text{syni} = 0.05\ \text{nS}$
  \item $k_\text{syn} = 0.25$, \quad $V_\text{th} = -52.0\ \text{mV}$
  \item $I_\text{app} = 3.0\ \text{pA}$ on $N_0$ --- below 2 pA: no oscillation; above 4 pA: locked above threshold
  \item $\beta = 0.2$, \quad $\tau_w = 200\ \text{ms}$
\end{itemize}

\textbf{Sweep:} \texttt{tests/sweep\_two\_EI.py} updated to remove SF and run a $G_\text{syne} \times G_\text{syni} \times I_\text{app}$ grid (0.01--1.0 nS, $I_\text{app}$ 2--5 pA). Confirmed oscillations exist in a region around the validated parameters.

\textbf{Note on sine-wave input:} Earlier tests used oscillatory sine-wave input ($f = 0.002$--$0.004\ \text{cyc/ms}$). Oscillations at $f = 0.004$ broke down because the input period (250 ms) is comparable to $\tau_w$ (200 ms), preventing full $w$ recovery between cycles. The constant-$I_\text{app}$ test is more informative for assessing intrinsic oscillatory capacity.

\textbf{Test file:} \texttt{tests/test\_two\_EI.py}

\section{Single-Segment Network Test}

\textbf{Objective:} Apply the EI pacemaker concept to the first segment (17 cells) of the full 102-cell network.

\subsection{Network Structure}

Each segment contains 17 cells across 9 classes:

\begin{center}
\begin{tabular}{@{}clll@{}}
  \toprule
  Class & Cells  & Count & Type \\
  \midrule
  1 & 0--1   & 2 & Interneuron \\
  2 & 2      & 1 & Interneuron \\
  3 & 3      & 1 & Interneuron \\
  4 & 4      & 1 & Interneuron \\
  5 & 5--6   & 2 & Motor neuron \\
  6 & 7--8   & 2 & Motor neuron \\
  7 & 9--10  & 2 & Motor neuron \\
  8 & 11--13 & 3 & Dorsal muscle (passive) \\
  9 & 14--16 & 3 & Ventral muscle (passive) \\
  \bottomrule
\end{tabular}
\end{center}

\subsection{Network Directionality}

Inter-segment connectivity is asymmetric:
\begin{itemize}
  \item \textbf{Forward (anterior $\to$ posterior):} Classes 1, 3, 6 from segment $n$ excite \textbf{class 4} in segment $n+1$ --- an active oscillator cell.
  \item \textbf{Backward (posterior $\to$ anterior):} Class 7 from segment $n$ excites class 4 in segment $n-1$, but classes 1 and 2 project onto \textbf{class 8} (passive muscle) --- a dead end for signal propagation.
\end{itemize}
This structural asymmetry biases signal propagation in the anterior-to-posterior direction, consistent with the forward locomotion travelling wave.

\subsection{Class-1 Connectivity (Within Segment)}

\begin{center}
\begin{tabular}{@{}cllll@{}}
  \toprule
  Cell & Excitatory OUT & Excitatory IN & Gap junctions \\
  \midrule
  0 & classes 2, 4, 5, 8 & class 3 & --- \\
  1 & classes 5, 8, 8    & class 3 & class 2 \\
  \bottomrule
\end{tabular}
\end{center}

Cell 0 is the more influential driver --- it connects to class 4, which is the recipient of the inter-segment forward wave.

\subsection{EI Pacemaker Design}

Since the segment does not self-oscillate with constant drive, two external ``head'' inhibitory neurons (cells 17--18) were added, forming EI loops with the two class-1 cells:
\begin{itemize}
  \item Cell 0 $\leftrightarrow$ Head neuron 17
  \item Cell 1 $\leftrightarrow$ Head neuron 18
\end{itemize}
$I_\text{app} = 3.0\ \text{pA}$ applied to cells 0 and 1 only. Head neurons receive no direct input and have no SF. Segment internal connectivity retains SF.

\textbf{Status:} EI loops oscillate as intended. Propagation into the rest of the segment through SF-gated synapses is still under investigation.

\textbf{Test file:} \texttt{tests/test\_single\_segment.py}

\section{Open Questions}

\begin{enumerate}
  \item Do the EI-driven class-1 oscillations propagate reliably to classes 2--7, or does SF in the segment connectivity attenuate them too strongly?
  \item Should SF be replaced by $w$ throughout the segment, or retained as a complementary mechanism?
  \item Once the single segment works, how should inter-segment coupling be reintroduced --- all six segments at once, or one at a time?
  \item Class 4 is missing its inter-segment excitatory drive in the single-segment test (classes 1, 3, 6 from the previous segment). This may need to be added as a synthetic input to get class 4 oscillating reliably.
\end{enumerate}

\end{document}
